% ID: SEL-D0020-OA-E01
% ESTADO: cotejada. Segunda pasada visual completa y borrador compilado revisados.
% Fuente: Selectividad/matematicas_aplicadas/reserva b/Examen-Matemáticas Aplicadas a las Ciencias Sociales.pdf
% SHA-256: 37584a9604930b74008a6bda88ef3d61fa926d5e6a540ab9013569258f772256
% Página: 1; opción A; ejercicio original 1.

\textbf{(2.5 puntos)} Una joyería elabora dos tipos de collares a partir de perlas blancas, grises y negras.

Para un collar de tipo A hacen falta 20 perlas blancas, 20 grises y 30 negras, mientras que para un collar del tipo B, 10 perlas blancas, 20 grises y 60 negras. Se dispone de un máximo de 900 perlas blancas y 1400 grises, mientras que es necesario que se utilicen al menos 1800 perlas negras.

Sabiendo que cada collar del tipo A le supone a la joyería un beneficio de 600 euros y cada collar del tipo B, 500 euros, calcule cuál debe ser la producción para obtener el máximo beneficio, así como a cuánto asciende el mismo. \textquestiondown{}Es posible fabricar 40 collares del tipo A y 20 del tipo B?
